\section{The Quiet Days of January and the Threshold of the Epiphany}\label{sec:jan-ferias}

Between the Octave Day on 1~January and First Vespers of the Epiphany on the
evening of 5~January lie four quiet calendar days. One of them---always a
Sunday when one falls there, or 2~January otherwise---belongs to the Most Holy
Name of Jesus (Section~\ref{sec:holy-name}). The remaining three days bring
Christmastide to a serene close through three distinct liturgical threads in
the 1962 Roman Rite: the ordinary weekdays of 2--5~January, the Saturday Office
and Mass of Our Lady when a Saturday falls in those days, and the commemoration
of Pope St.~Telesphorus on 5~January.

\subsection{The Weekdays of 2--5 January: Lingering in the Octave's Joy}

In pre-1955 Roman books, 2, 3, and 4~January were the octave days of
St.~Stephen, St.~John, and the Holy Innocents, and 5~January was the Vigil of
the Epiphany. Under the 1960 \textit{Codex Rubricarum}, all four of those
observances disappeared (General Rubrics~30--32 and~64), leaving the non-feast
weekdays from Monday through Friday in 2--5~January as fourth-class ferias of
Christmastide (RG~91, row~28).

Ordinarily a fourth-class feria repeats the preceding Sunday's Mass without the
\textit{Gloria} (RGMR~299). Here, however, the 1962 Missal (p.~34) makes a
deliberate seasonal exception: on the weekdays from 2 to 5~January, the priest
celebrates the \textbf{Mass of 1~January} (\textit{Puer natus est}, Collect
\textit{Deus, qui salutis aeternae}, Epistle Tit~2:11--15, Gospel Luke~2:21) in
white vestments, \textbf{with the \textit{Gloria} and the Nativity Preface},
yet \textbf{without the Creed and without the proper \textit{Communicantes}}.
The liturgy lowers its voice just enough to mark that the eight-day Octave has
closed, while keeping the joyful hymn of the angels and the Marian prayer of
1~January on the altar through the fifth day of the year.

\subsection{Our Lady on Saturday in 2--5 January}

Whenever 25~December falls on Tuesday, Wednesday, Thursday, or Friday---fifty-seven
seasons in every century---one of the four days from 2 to 5~January falls on a
\textbf{Saturday} (2~January when Christmas is Friday; 3~January when Christmas
is Thursday; 4~January when Christmas is Wednesday; and 5~January when
Christmas is Tuesday). Because the Most Holy Name of Jesus can never land on a
Saturday (its 2~January fallback occurs only on Monday, Tuesday, or Wednesday),
every Saturday inside 2--5~January would otherwise be an ordinary fourth-class
feria.

Under General Rubric~78 (\textit{fit de sancta Maria in sabbato}) and General
Rubric~91 (row~27 outranking row~28), the \textbf{Office of Our Lady on
Saturday} is obligatory on every fourth-class Saturday, and by Missal
Rubrics~270 and~515~b the Mass agrees with the Office of the day. On a Saturday
in 2--5~January, the 1962 Missal therefore appoints not the 1~January Mass, but
the second seasonal Mass of Our Lady on Saturday (\textit{A Nativitate Domini
usque ad Purificationem}, pp.~[44]--[45]): \textbf{\textit{Vultum tuum}}.

The Mass of \textit{Vultum tuum} is one of the hidden jewels of Christmastide.
Its Introit draws from the bridal Psalm~44~[45]:13: \textit{Vultum tuum
deprecabuntur omnes divites plebis}, ``All the rich among the people shall
entreat thy countenance.'' Its Collect is the same prayer of 1~January,
\textit{Deus, qui salutis aeternae}. Its Epistle and Gospel return to the
Christmas Dawn Mass---Titus~3:4--7 and the shepherds hurrying to Bethlehem in
Luke~2:15--20, where St.~Ambrose (\textit{Exposition of Luke}, Book~II,
section~53) bids us learn ``the modesty of the Holy Virgin, that modesty of
speech as of body, whereby she laid up in her heart the evidences of her
faith''---while its Alleluia sings the Christmastide versicle of \textit{Alma
Redemptoris Mater}: \textit{Post partum, Virgo, inviolata permansisti: Dei
Genetrix, intercede pro nobis}, ``After childbirth thou didst still remain an
inviolate virgin: O Mother of God, intercede for us.'' The Offertory proclaims
\textit{Felix namque es, sacra Virgo Maria}, ``For thou art happy, O sacred
Virgin Mary, and most worthy of all praise, since out of thee hath arisen the
sun of justice, Christ our God,'' and the Communion echoes the Beatitude of
Luke~11:27: \textit{Beata viscera Mariae Virginis, quae portaverunt aeterni
Patris Filium}, ``Blessed is the womb of the Virgin Mary, which bore the Son of
the eternal Father.''

One rubric detail is worth noting cleanly: \textit{Vultum tuum} takes the
\textit{Gloria}, no Creed, the common \textit{Communicantes}, and the
\textbf{Preface of the Blessed Virgin Mary} (\textit{Et te in veneratione},
explicitly appointed on p.~[45] of the 1962 Missal), rather than the Nativity
Preface. Within the Christmas Octave (26--31~December), the Nativity Preface
overrides proper prefaces of saints; on 2--5~January, it is used only as a
seasonal preface without that override, leaving Our Lady's Saturday Mass free
to sing her own proper preface.

\subsection{5 January: Pope St.~Telesphorus, Martyr}

On 5~January---whenever it falls on Monday through Saturday (eighty-seven years
in a century, including Saturday 5~January when \textit{Vultum tuum} is
celebrated)---Low Mass adds an ordinary commemoration of \textbf{Pope
St.~Telesphorus, Martyr} (omitted in non-conventual Sung Mass, and omitted
completely when 5~January is Sunday and the Most Holy Name of Jesus is
celebrated). His orations come from the Common of Supreme Pontiffs with the
martyr's title inserted: \textit{Gregem tuum, Pastor aeterne... per beatum
Telesphorum Martyrem tuum atque Summum Pontificem}, asking the Eternal
Shepherd to guard his flock through his Martyr and Supreme Pontiff.

The \textit{Roman Martyrology} for 5~January honors St.~Telesphorus---a Greek
by birth who governed the Church of Rome in the early second century
(c.~125--136)---for his glorious martyrdom under Emperor Antoninus Pius, an
ancient martyrdom already attested in the second century by St.~Irenaeus.
Medieval liturgists such as Sicard of Cremona (\textit{Mitrale}, following the
sixth-century \textit{Liber Pontificalis}) went further and preserved a
tradition crediting Pope St.~Telesphorus with instituting the Christmas
Midnight Mass and the singing of the \textit{Gloria in excelsis} on Christmas
night. While modern historians treat that liturgical attribution strictly as a
medieval tradition rather than verified second-century fact, St.~Telesphorus's
place on 5~January provides a fitting Roman close to the twelve days begun at
Midnight Mass.

\subsection{None of 5 January and the Turning Toward the Star}

When the Church finishes \textbf{None on 5~January}, Christmastide
(\textit{Tempus Nativitatis}, General Rubric~72~a) reaches its exact canonical
end. That same evening, \textbf{First Vespers of the Epiphany} opens
\textbf{Epiphanytide} (\textit{Tempus Epiphaniae}, General Rubric~72~b),
running from the evening of 5~January through 13~January inclusive.

Across these twelve days the liturgy has stayed close to the hidden stable of
Bethlehem: the silence of midnight, the shepherds and angels, the blood of
Stephen and the Innocents, the contemplative gaze of John, Simeon and Anna in
the Temple, the eighth-day Circumcision and the Virgin Mother on 1~January, the
sweetness of the Name of Jesus, and Mary's quiet Saturday beside the crib. Now,
as dusk gathers on 5~January, the doors of the stable swing wide toward the
nations: the star shines over the East, the Magi bring gold, frankincense, and
myrrh, the waters of the Jordan wait to be sanctified by the Lord's Baptism,
and the wine of Cana stands ready. In the Roman liturgy, the Epiphany is not
the end of Christmas; it is the unveiling to all the earth of the Child whom
the shepherds adored in the dark.
